\scanpage{35}
Définir $K^\bullet_\Phi(X)$\quad ($X$ préschéma.).

Définir un substitut de l'anneau de Chow
\[
\unsure{A}^\bullet_\Phi(X)
\]
(qui doit être un \inserted{pseudo}anneau \unsure{comm.} \inserted{(pas néc d'élément $1$) et une \unsure{flèche} vers $A^\bullet(X)$, l'anneau de Chow ``ordinaire'' \unsure{whitney th.\ + unaires|whitney th.\ + unitaires}]} et des classes de Chern
\[
c^i_\Phi:K^\bullet_\Phi(X)\mathbin{\text{---}}A_\Phi(X)\qquad\boxed{i>0}
\]
qui doit satisfaire la condition de multiplicativité ordinaire interne
\[
\unsure{c^n_\Phi}(x+y)=c^n_\Phi(x)+c^n_\Phi(y)+\sum_{p+q=n\atop p,q>0}c^p_\Phi(x)c^q_\Phi(y)
\]
$(x,y\in K^\bullet_\Phi(X))$, et \deleted{\unsure{pourront}} être compatibles avec \deleted{des} les hom \deleted{\unsure{can}} can.
\diagram{\begin{tikzcd}
K^\bullet_\Phi(X)\arrow[r,"i"]\arrow[d,dash,"c_\Phi"'] & K^\bullet(X)\arrow[d,"c"]\\
A^\bullet_\Phi(X)\arrow[r,"j"] & A^\bullet(X)
\end{tikzcd}}
où la dernière est un hom multiplicatif (canonique).

Cas \inserted{particuliers} à \deleted{\unsure{distinguer}} étudier

1) $X$ est lisse qu.\ pr. sur un corps $k$, \unsure{on prend la} définition habituelle de l'équivalence rationnelle, mais en prenant tjrs des cycles sur $X$ à support $\unsure{\subset|\in}\Phi$.

2) $X$ général: on essaiera de définir $A^\bullet_\Phi(X)$ via une filtration sur le $\lambda$-pseudo-anneau $K^\bullet_\Phi(X)$ ??

3) Classes de Chern \unsure{cohomologiques|homologiques} ??
\margin{34}
